% ======================================================================
% INF5163 – Projet de session 2 – Rapport principal
% Compilation : « make rapport » depuis la racine → build/rapport.pdf
% Convention : une phrase par ligne dans tous les fichiers .tex.
% ======================================================================
\documentclass[11pt]{article}
% ======================================================================
% Préambule du rapport (partagé avec rapport/entente-autonome.tex).
% Caractères grecs : en mode mathématique ($\varepsilon$, $\delta$) avec pdflatex.
% Chemins relatifs à la RACINE du dépôt : compiler avec « make » ou
% latexmk lancé depuis la racine (Overleaf compile aussi depuis la racine).
% ======================================================================

% --- Langue et encodage (pdflatex) ---------------------------------------------
\usepackage[T1]{fontenc}
\usepackage{lmodern}
\usepackage[french]{babel}
\usepackage[autostyle=true]{csquotes}   % \enquote{…} → « … »

% --- Mise en page ----------------------------------------------------------------
\usepackage[letterpaper,margin=2.5cm]{geometry}
\usepackage{microtype}
\usepackage{setspace}
\onehalfspacing % TODO (équipe) : confirmer l'interligne attendu (simple / 1,5).

% --- Tableaux, figures, couleurs --------------------------------------------------
\usepackage{booktabs}
\usepackage{tabularx}
\usepackage{array}
\usepackage{graphicx}
\graphicspath{{shared/figures/}}
\usepackage[table]{xcolor}
\usepackage{enumitem}
\usepackage{amsmath}               % paramètres de DP : ε, δ, etc.
\usepackage[titletoc,title]{appendix} % « Annexe A », « Annexe B »
\usepackage{refcount}              % contrôle de la longueur (rapport/main.tex)

% --- Bibliographie ----------------------------------------------------------------
% Style IEEE (paquet biblatex-ieee). Pour passer en APA : style=apa.
\usepackage[backend=biber,style=ieee]{biblatex}
\addbibresource{shared/references.bib}

% --- Macros partagées -------------------------------------------------------------
% ======================================================================
% Macros partagées par le rapport, l'entente autonome et la présentation.
% Ce fichier est chargé (\input) dans les préambules : ne rien y composer.
% ======================================================================

% --- Noms propres et cadres ------------------------------------------------
% Utilisation : \fedesante{}, \pcap{}, \eptc{} (les accolades évitent
% que l'espace suivante soit avalée).
\newcommand{\fedesante}{FédéSanté}
% PCAP est une marque déposée du Centre de gouvernance de l'information des
% Premières Nations (CGIPN) ; le symbole ® est demandé par le CGIPN.
% TODO : vérifier la forme d'usage recommandée par le CGIPN.
\newcommand{\pcap}{PCAP\textsuperscript{\textregistered}}
% TODO : vérifier la version de l'EPTC 2 en vigueur à la date de remise
% et ajuster l'entrée correspondante dans shared/references.bib.
\newcommand{\eptc}{EPTC~2}
\newcommand{\loivingtcinq}{Loi~25}

% --- Espace réservé visible ----------------------------------------------
% Marque dans le PDF ce qui reste à rédiger. Aucun \arediger ne doit
% subsister dans la version remise (vérifier avec : grep -rn arediger).
\newcommand{\arediger}[1]{{\color{gray}\itshape[À rédiger : #1]}}

% --- Sigles et acronymes (paquet acronym) ----------------------------------
% \ac{FL} : forme longue + sigle au premier emploi, sigle ensuite.
% \acs{FL} : sigle seul ; \acl{FL} : forme longue seule.
\usepackage{acronym}
\acrodef{FL}[FL]{apprentissage fédéré (\emph{federated learning})}
\acrodef{DP}[DP]{confidentialité différentielle (\emph{differential privacy})}
\acrodef{FedAvg}[FedAvg]{moyenne fédérée (\emph{Federated Averaging})}
\acrodef{CER}[CER]{comité d'éthique de la recherche}
\acrodef{PCAP}[PCAP]{propriété, contrôle, accès et possession}


% --- Liens et renvois (hyperref, puis cleveref, en dernier) -----------------------------
\usepackage[hidelinks,pdfusetitle]{hyperref}
\hypersetup{
  pdftitle={Quand le modèle voyage : éthique, propriété et souveraineté dans un consortium d'apprentissage fédéré en santé},
  pdfsubject={INF5163 -- Projet de session 2},
}
\usepackage[french,nameinlink]{cleveref}

% --- Articles numérotés de l'entente (section 8) --------------------------------
% \articleentente{Titre} produit « Article N – Titre » et l'ajoute à la table
% des matières. \label{art:...} juste après permet \ref{art:...}.
\newcounter{articleentente}
\newcommand{\articleentente}[1]{%
  \refstepcounter{articleentente}%
  \subsection*{Article~\thearticleentente{} -- #1}%
  \addcontentsline{toc}{subsection}{Article~\thearticleentente{} -- #1}}

% Vrai uniquement dans rapport/entente-autonome.tex.
\newif\ifententeautonome
\ententeautonomefalse


% --- Contrôle de la longueur du corps (15 à 20 pages) -----------------------------
% Le corps va de \label{debut-corps} (début de l'introduction) à
% \label{fin-corps} (fin de la conclusion), hors références et annexes.
% Le nombre de pages est écrit dans le log ; « make pages » l'affiche.
\AtEndDocument{%
  \IfRefUndefinedExpandable{fin-corps}{%
    \typeout{FEDESANTE-PAGES-CORPS: inconnu (recompiler)}%
  }{%
    \typeout{FEDESANTE-PAGES-CORPS: \number\numexpr\getpagerefnumber{fin-corps}-\getpagerefnumber{debut-corps}+1\relax}%
  }%
}

\begin{document}

% --- Pages liminaires -------------------------------------------------------------------
\pagenumbering{roman}
% ======================================================================
% Page titre. TODO : remplacer les espaces réservés par les noms réels
% (ordre alphabétique de nom de famille, par exemple) et retirer la
% ligne 8 si l'équipe compte 7 membres.
% ======================================================================
\begin{titlepage}
  \centering
  {\large Université du Québec en Outaouais (UQO)\par}
  {\large Département d'informatique et d'ingénierie\par} % TODO : vérifier l'intitulé exact du département.
  \vspace{2.5cm}

  {\LARGE\bfseries Quand le modèle voyage\par}
  \vspace{0.4cm}
  {\Large Éthique, propriété et souveraineté dans un consortium d'apprentissage fédéré en santé\par}
  \vspace{2cm}

  {\large INF5163 -- Méthodologie de recherche en informatique\par}
  {\large Projet de session 2 (examen final)\par}
  \vspace{0.6cm}
  {Présenté à N'dah Daniel Yapi, Ph.D.\par}
  \vspace{1.5cm}

  {\bfseries Membres de l'équipe\par}
  \vspace{0.3cm}
  \begin{tabular}{ll}
    \toprule
    Nom & Code permanent \\ % TODO : confirmer l'identifiant demandé.
    \midrule
    {[Membre 1]} & {[XXXX00000000]} \\
    {[Membre 2]} & {[XXXX00000000]} \\
    {[Membre 3]} & {[XXXX00000000]} \\
    {[Membre 4]} & {[XXXX00000000]} \\
    {[Membre 5]} & {[XXXX00000000]} \\
    {[Membre 6]} & {[XXXX00000000]} \\
    {[Membre 7]} & {[XXXX00000000]} \\
    {[Membre 8]} & {[XXXX00000000]} \\
    \bottomrule
  \end{tabular}

  \vfill
  {Automne 2026\par}
  {Remis le [date de remise]\par} % TODO : date de remise sur Moodle.
\end{titlepage}

\tableofcontents
\clearpage

% --- Corps du rapport (15 à 20 pages) --------------------------------------------------
\pagenumbering{arabic}
% Numérotation alignée sur l'énoncé et la grille (1 = page titre) : l'introduction
% porte le numéro 2, l'analyse éthique le numéro 5, etc. Retirer la ligne
% suivante pour numéroter à partir de 1.
\setcounter{section}{1}
\label{debut-corps}
% !TEX root = ../main.tex
% ======================================================================
% Section 2 – Introduction et justification du sujet
% Grille : « Cadrage & problématique » – 6 pts (partagés avec la section 3).
% Longueur indicative : ~1,5 à 2 pages.
% ======================================================================
% TODO : situer le sujet dans la recherche multicentrique en santé ou dans la protection de la vie privée.
% TODO : expliquer les limites des approches centralisées (regroupement des dossiers patients dans un entrepôt unique).
% TODO : présenter brièvement \fedesante{} (mise en situation hypothétique) et la promesse « les données ne quittent jamais l'établissement ».
% TODO : annoncer le plan du rapport (renvois avec \cref{sec:...}).

\subsection{Recherche multicentrique en santé et protection de la vie privée}
\label{sec:intro-contexte}
\arediger{contexte}

\subsection{Limites des approches centralisées}
\label{sec:intro-limites}
\arediger{limites}

\subsection{Mise en situation et plan du rapport}
\label{sec:intro-plan}
\arediger{mise en situation et plan}

% !TEX root = ../main.tex
% ======================================================================
% Section 3 – Rappel de la problématique et des objectifs
% Grille : « Cadrage & problématique » – 6 pts (partagés avec la section 2).
% Longueur indicative : ~1 à 1,5 page.
% ======================================================================
% TODO : ce que \fedesante{} visait à résoudre (ex. « entraîner un modèle performant sans centraliser les dossiers patients »).
% TODO : formuler des objectifs SMART (spécifiques, mesurables, atteignables, réalistes, temporellement définis).
% TODO : énoncer le mandat de l'équipe : analyser de manière critique si la promesse suffit, et quels enjeux éthiques, légaux et de gouvernance subsistent.

\subsection{Problématique visée par \fedesante{}}
\label{sec:problematique-visee}
\arediger{problématique}

\subsection{Objectifs SMART}
\label{sec:objectifs-smart}
\arediger{objectifs}
% Suggestion : un tableau booktabs Objectif | Indicateur | Cible | Échéance.

\subsection{Mandat de l'analyse critique}
\label{sec:mandat}
\arediger{mandat}

% !TEX root = ../main.tex
% ======================================================================
% Section 4 – Rappel du schéma méthodologique
% Responsable : Frank Onana Tsanga
% Grille : pas de critère propre ; soutient « Cadrage & problématique »
% et sert de base technique aux sections 5 et 6.
% Longueur indicative : ~1,5 à 2 pages.
% ======================================================================
% TODO : décrire l'architecture fédérée (12 établissements : hôpitaux universitaires, cliniques régionales, 3 centres de santé de Premières Nations ; serveur d'agrégation).
% TODO : présenter \ac{FedAvg} pondéré par la taille des sites (citer \cite{mcmahan2017fedavg}).
% TODO : préciser les paramètres de \ac{DP} : ε, δ, écrêtage (clipping), bruit ; niveau où la DP est appliquée (local ou central).
% TODO : figure de l'architecture dans shared/figures/ (\includegraphics).

\subsection{Architecture fédérée}
\label{sec:architecture}
\arediger{architecture}

\subsection{Agrégation : FedAvg pondéré par la taille des sites}
\label{sec:fedavg}
\arediger{agrégation}

\subsection{Paramètres de confidentialité différentielle}
\label{sec:parametres-dp}
\arediger{paramètres $\varepsilon$, $\delta$, écrêtage, bruit}

% !TEX root = ../main.tex
% ======================================================================
% Section 5 – Analyse d'intégrité, éthique et ÉDI
% Grille : « Analyse éthique & ÉDI » – 16 pts (section prioritaire).
% Longueur indicative : ~4 à 5 pages.
% Cadres obligatoires : \eptc{} (dont le chapitre 9) et la \loivingtcinq{}.
% ======================================================================
% FIL CONDUCTEUR (toute la section) : analyse critique de la promesse
% « les données ne quittent jamais l'établissement ».

\subsection{Fil conducteur : la promesse « les données ne quittent jamais l'établissement »}
\label{sec:promesse}
% TODO : ce que la promesse garantit réellement et ce qu'elle ne couvre pas.
\arediger{analyse critique de la promesse}

\subsection{Intégrité et éthique}
\label{sec:integrite-ethique}

\subsubsection{Risques résiduels}
\label{sec:risques-residuels}
% TODO : quels risques résiduels subsistent, dont l'inférence d'appartenance et l'inversion de modèle à partir des mises à jour ?
% TODO : la confidentialité différentielle couvre-t-elle ces risques ? À quel prix (compromis ε / performance) ?
\arediger{risques résiduels}

\subsubsection{Choix du modèle de consentement}
\label{sec:consentement}
% TODO : quel modèle de consentement choisir : large, spécifique ou dynamique ? Justifier au regard de l'EPTC 2 et de la Loi 25.
\arediger{consentement}

\subsubsection{Usages secondaires du modèle}
\label{sec:usages-secondaires}
% TODO : quels usages secondaires du modèle sont envisageables, et lesquels doivent être encadrés ou interdits ?
\arediger{usages secondaires}

\subsubsection{Proposition d'un protocole de gouvernance}
\label{sec:protocole-gouvernance}
% TODO : proposer un protocole de gouvernance (instances, rôles, décisions, rôle des CER et des instances autochtones).
\arediger{protocole de gouvernance}

\subsection{Équité, diversité et inclusion (ÉDI)}
\label{sec:edi}

\subsubsection{Agrégation pondérée et petits sites}
\label{sec:edi-petits-sites}
% TODO : l'agrégation pondérée par la taille marginalise-t-elle les petits sites ?
\arediger{petits sites}

\subsubsection{Performance par sous-groupe}
\label{sec:edi-sous-groupes}
% TODO : le modèle performe-t-il également pour chaque sous-groupe ? Comment le mesurer ?
\arediger{sous-groupes}

\subsubsection{Prévention de la stigmatisation}
\label{sec:edi-stigmatisation}
% TODO : comment éviter la stigmatisation des communautés plus touchées par le diabète ?
\arediger{stigmatisation}

% !TEX root = ../main.tex
% ======================================================================
% Section 6 – Gestion, responsabilité et propriété intellectuelle
% Responsable : Alex D'Andrea
% Grille : « PI, gestion des données & responsabilité » – 14 pts (section prioritaire).
% Longueur indicative : ~3,5 à 4 pages.
% ======================================================================
\subsection{Propriété intellectuelle et responsabilité}
\label{sec:pi-responsabilite}
Cette section traite de la propriété intellectuelle et des responsabilités de chaque acteur.
Dans un modèle fédéré, la contribution est partagée entre plusieurs établissements, ce qui rend la question de la \ac{PI} plus complexe.

\subsubsection{Propriété du modèle global}
\label{sec:propriete-modele}
Le code source, le protocole de traitement des données et le nom \fedesante{} appartiennent à l'équipe qui les a conçus.
Le modèle, en soi, pose une autre question.
Sa valeur réside dans ses poids : sans eux, l'architecture est inutile.
Or, ces poids ont été ajustés à partir des données des 12 établissements participants.
\arediger{confirmer l'algo d'agrégation dans section metho}
Avec \ac{FedAvg} pondéré, les grands sites pèsent davantage dans l'agrégation, mais chaque contribution reste indispensable au modèle final.
Les auteurs proposent donc une copropriété du modèle global, gérée par un consortium.
Chaque établissement y disposera d'une voix, quelle que soit sa taille, afin d'assurer l'équité dans la gestion du modèle.
Cette copropriété ne suffit toutefois pas pour les trois centres de santé des Premières Nations, dont les données relèvent d'une gouvernance propre, abordée à la \cref{sec:pcap-parametres}.


\subsubsection{Principes PCAP appliqués aux paramètres du modèle}
\label{sec:pcap-parametres}
% TODO : comment appliquer les principes \pcap{} (propriété, contrôle, accès, possession) à des paramètres de modèle plutôt qu'à des données brutes ? (citer \cite{pcap})
\arediger{PCAP et paramètres}

\subsubsection{Retrait d'un site et désapprentissage}
\label{sec:retrait-desapprentissage}
% TODO : un site peut-il se retirer et exiger le « désapprentissage » (machine unlearning) de sa contribution ? Faisabilité technique et engagements contractuels.
\arediger{retrait et désapprentissage}

\subsubsection{Protocole d'agrégation : brevet ou code ouvert ?}
\label{sec:brevet-code-ouvert}
% TODO : le protocole d'agrégation doit-il être breveté ou publié en code ouvert ?
\arediger{brevet ou code ouvert}

\subsection{Documentation}
\label{sec:documentation}

\subsubsection{Fiche de modèle (\emph{model card})}
\label{sec:model-card}
% TODO : contenu de la fiche de modèle.
\arediger{fiche de modèle}

\subsubsection{Fiche de données par site (\emph{datasheet})}
\label{sec:datasheet}
% TODO : contenu de la fiche de données par site.
\arediger{fiche de données}

\subsubsection{Registre des participants}
\label{sec:registre}
% TODO : contenu et tenue du registre des participants.
\arediger{registre}

\subsubsection{Règles d'accès au modèle pour les chercheurs externes}
\label{sec:acces-externe}
% TODO : règles d'accès au modèle pour les chercheurs externes (académiques ou privés).
\arediger{règles d'accès}

% !TEX root = ../main.tex
% ======================================================================
% Section 7 – Valorisation et diffusion
% Responsable : Stive Roussel Tankio Djiokap
% Grille : « Valorisation & diffusion » – 8 pts (section prioritaire).
% Longueur indicative : ~2 à 2,5 pages.
% ======================================================================
\subsection{Valorisation}
\label{sec:valorisation-modele}
% TODO : plateforme publique de recherche ET/OU licence aux fabricants d'appareils médicaux ?
% TODO : analyser la tension entre ces deux voies.
% TODO : partage des bénéfices avec les communautés participantes.
\arediger{valorisation}

\subsection{Diffusion}
\label{sec:diffusion}

\subsubsection{Comités d'éthique de la recherche et instances de gouvernance autochtones}
\label{sec:diffusion-cer}
% TODO : public 1 — les comités d'éthique de la recherche (\acs{CER}) et les instances de gouvernance autochtones (\eptc{}, chapitre 9).
\arediger{diffusion auprès des CER et des instances autochtones}

\subsubsection{Patients et vulgarisation}
\label{sec:diffusion-patients}
% TODO : public 2 — les patients, avec un volet vulgarisation.
% TODO : comment les rassurer sur la sécurité et la pertinence du projet ?
\arediger{diffusion auprès des patients}

% ======================================================================
% Section 8 – Livrable spécifique : entente de gouvernance des données simplifiée
% Grille : « Livrable spécifique » – 10 pts (section prioritaire).
% Longueur : 2 à 3 pages (vérifier avec « make entente » → build/entente.pdf).
%
% Ce fichier est inclus DANS LE RAPPORT et DANS rapport/entente-autonome.tex :
% ne rien y écrire qui n'ait de sens que dans le rapport, sauf dans le
% bloc \ifententeautonome\else … \fi ci-dessous.
%
% Exigences : entente entre le consortium et UNE communauté participante,
% claire, applicable et cohérente avec les principes PCAP. Articles
% numérotés avec \articleentente{Titre}.
% ======================================================================
\setcounter{articleentente}{0}

\ifententeautonome\else
\section{Livrable spécifique : entente de gouvernance des données simplifiée}
\label{sec:entente}
% TODO : un court paragraphe de présentation (rapport seulement) : rôle de l'entente, lien avec les sections 5 et 6.
\arediger{présentation de l'entente}
\fi

\articleentente{Parties}
\label{art:parties}
% TODO : identifier le consortium \fedesante{} (représentant) et la communauté participante (instance de gouvernance, centre de santé).
\arediger{parties}

\articleentente{Objet}
\label{art:objet}
% TODO : objet de l'entente (participation à l'entraînement fédéré, portée, exclusions).
\arediger{objet}

\articleentente{Définitions}
\label{art:definitions}
% TODO : définir au minimum : données locales, mise à jour du modèle, modèle global, \acs{DP}, désapprentissage, publication, bénéfices.
\arediger{définitions}

\articleentente{Droits PCAP de la communauté}
\label{art:droits-pcap}
% TODO : traduire chacun des quatre principes (propriété, contrôle, accès, possession) en droits concrets et vérifiables.
\arediger{droits PCAP}

\articleentente{Conditions de retrait et désapprentissage}
\label{art:retrait}
% TODO : préavis, effets du retrait, engagement de désapprentissage (ou de réentraînement sans la contribution), délais, preuve fournie.
\arediger{retrait et désapprentissage}

\articleentente{Approbation des publications}
\label{art:publications}
% TODO : processus d'approbation préalable des publications et présentations par la communauté ; délais ; droit de refus.
\arediger{approbation des publications}

\articleentente{Partage des bénéfices}
\label{art:benefices}
% TODO : bénéfices (accès au modèle, retombées financières, renforcement des capacités, reconnaissance).
\arediger{partage des bénéfices}

\articleentente{Durée}
\label{art:duree}
% TODO : durée, renouvellement, sort du modèle et des documents à l'échéance.
\arediger{durée}

\articleentente{Résolution des différends}
\label{art:differends}
% TODO : étapes (discussion, médiation, instance désignée) ; respect des modes de gouvernance de la communauté.
\arediger{résolution des différends}

\articleentente{Signatures}
\label{art:signatures}
% TODO : vérifier les titres des signataires.
\vspace{0.5em}
\noindent
\begin{tabularx}{\linewidth}{@{}XX@{}}
Pour le consortium \fedesante{} & Pour la communauté participante \\[2.5em]
\rule{0.9\linewidth}{0.4pt} & \rule{0.9\linewidth}{0.4pt} \\
Nom, titre : [\ldots] & Nom, titre : [\ldots] \\
Date : [\ldots] & Date : [\ldots] \\
\end{tabularx}

% !TEX root = ../main.tex
% ======================================================================
% Section 9 – Conclusion
% Responsable : ___ (non attribuée : à décider en équipe)
% Grille : pas de critère propre ; contribue à « Cadrage & problématique »
% (cohérence avec les objectifs) et à « Rédaction & normes ».
% Longueur indicative : ~1 page.
% ======================================================================
% TODO : répondre au mandat : la promesse « les données ne quittent jamais l'établissement » suffit-elle ?
% TODO : synthèse des recommandations des sections 5 à 8.

\subsection{Synthèse}
\label{sec:synthese}
\arediger{synthèse}

\subsection{Limites de l'analyse}
\label{sec:limites}
% TODO : limites de l'analyse (cas hypothétique, cadres juridiques vérifiés à une date donnée, etc.).
\arediger{limites}

\subsection{Risques résiduels}
\label{sec:risques-conclusion}
% TODO : risques résiduels qui subsistent malgré les mesures proposées.
\arediger{risques résiduels}

% Fin du corps du rapport : NE PAS DÉPLACER (contrôle des 15 à 20 pages).
\label{fin-corps}
 % contient \label{fin-corps}

% --- Références (hors limite de pages) -----------------------------------------------------
\clearpage
% TODO : retirer ce \nocite quand les sources d'exemple seront citées
% dans le texte avec \cite{...} (sinon elles n'apparaissent pas).
\nocite{eptc2,loi25,pcap,mcmahan2017fedavg}
\printbibliography[heading=bibintoc,title={Références}]

% --- Annexes (obligatoires, hors limite de pages) ---------------------------------------
\clearpage
\begin{appendices}
% ======================================================================
% Annexe A – Contributions individuelles (obligatoire, hors limite de pages)
% TODO : remplir une ligne par membre ; chaque membre signe.
% ======================================================================
\section{Contributions individuelles}
\label{ann:contributions}

{\renewcommand{\arraystretch}{1.3}
\begin{tabularx}{\linewidth}{@{}p{3.5cm}Xp{2cm}p{3.2cm}@{}}
\toprule
\textbf{Nom} & \textbf{Tâches principales (2-3 lignes)} & \textbf{Heures estimées} & \textbf{Signature} \\
\midrule
\rule{0pt}{3.5em} & & & \\ \midrule
\rule{0pt}{3.5em} & & & \\ \midrule
\rule{0pt}{3.5em} & & & \\ \midrule
\rule{0pt}{3.5em} & & & \\ \midrule
\rule{0pt}{3.5em} & & & \\ \midrule
\rule{0pt}{3.5em} & & & \\ \midrule
\rule{0pt}{3.5em} & & & \\ \midrule
\rule{0pt}{3.5em} & & & \\
\bottomrule
\end{tabularx}}

\medskip
\noindent\textbf{Note :} Toute divergence majeure de contribution doit être signalée à l'enseignant avant la remise.

\clearpage
% ======================================================================
% Annexe B – Déclaration d'intégrité et d'originalité (obligatoire)
% TODO : chaque membre relit la version finale, puis signe.
% ======================================================================
\section{Déclaration d'intégrité et d'originalité}
\label{ann:integrite}

Les membres de l'équipe soussignés déclarent que :
\begin{enumerate}[label=\arabic*.]
  \item le travail est original et n'est pas un plagiat, y compris de travaux d'équipes des sessions précédentes ;
  \item toutes les sources sont référencées ;
  \item l'usage de l'IA est décrit ;
  \item chaque membre a relu et approuvé la version finale.
\end{enumerate}

{\renewcommand{\arraystretch}{1.3}
\begin{tabularx}{\linewidth}{@{}p{4.5cm}Xp{4cm}@{}}
\toprule
\textbf{Nom} & \textbf{Courriel UQO} & \textbf{Signature} \\
\midrule
\rule{0pt}{2em} & & \\ \midrule
\rule{0pt}{2em} & & \\ \midrule
\rule{0pt}{2em} & & \\ \midrule
\rule{0pt}{2em} & & \\ \midrule
\rule{0pt}{2em} & & \\ \midrule
\rule{0pt}{2em} & & \\ \midrule
\rule{0pt}{2em} & & \\ \midrule
\rule{0pt}{2em} & & \\
\bottomrule
\end{tabularx}}

\subsection*{Déclaration d'usage de l'IA générative}
\addcontentsline{toc}{subsection}{Déclaration d'usage de l'IA générative}
\label{ann:ia}
% TODO : consigner AU FUR ET À MESURE chaque usage d'un outil d'IA générative
% (outil et version, portion précise du travail, raison). Ajouter des lignes au besoin.
% La première ligne est déjà remplie : la structure du dépôt a été générée avec Claude Code.

{\renewcommand{\arraystretch}{1.3}
\begin{tabularx}{\linewidth}{@{}p{3.5cm}XX@{}}
\toprule
\textbf{Outil} & \textbf{Portion du travail} & \textbf{Raison} \\
\midrule
Claude (Anthropic), via Claude Code & Structure initiale du dépôt : gabarits LaTeX vides, Makefile, intégration continue, fichiers de suivi (aucun contenu rédigé du rapport) & [À compléter par l'équipe] \\ \midrule
\rule{0pt}{2em} & & \\ \midrule
\rule{0pt}{2em} & & \\ \midrule
\rule{0pt}{2em} & & \\
\bottomrule
\end{tabularx}}

\end{appendices}

\end{document}
